% Section 3. The trial field built from candidate data.
% Rewritten 2026-11-08: the field is defined from cited components only
% (Barnett--Greengard cell fields, Botten et al. Bloch modes); solver-internal
% matrices are described in words. Numerical content unchanged.
\section{The trial field}\label{sec:field}

The bound of \cref{sec:bound} applies to one explicit field $\widehat u$ on the
strip, built from finitely many numbers: the \emph{candidate data}. This section
defines $\widehat u$ from these data. How the data are computed does not enter
the bound: whatever they are, $\widehat u$ solves the Helmholtz equation in
every homogeneous piece, and every inaccuracy of the data shows up as a jump of
$\widehat u$ across an interface.

We separate \emph{definition parameters} (density grids, Rayleigh truncation,
proxies, candidate data) from \emph{evaluation parameters} (target quadrature,
number of explicitly summed cells, evaluation precision). Once the former are
fixed, $\widehat u$ is a definite function; the latter only affect how
accurately the quantities of \cref{sec:bound} are bounded. Throughout,
$\widehat k>0$ is fixed, $\mu=\widehat k^2$, and $\kappa_{\mathrm e}=\widehat
k\rho_{\mathrm e}^{1/2}$.

\subsection{The field in one cell}\label{sec:layer}
\emph{Rayleigh term.} Fix $M\ge0$, let $K=2M+1$ and
$\gamma_m=(\kappa_{\mathrm e}^2-\beta_m^2)^{1/2}$ with $\Im\gamma_m\ge0$. For a cell
$(x_l,x_r)\times I_d$ and amplitudes $a,b\in\CC^K$, indexed by $|m|\le M$, let
\begin{equation}\label{eq:rayleigh}
  \mathcal R[a,b](x,y)=d^{-1/2}\sum_{|m|\le M}
  \bigl(a_m\ee^{\ii\gamma_m(x-x_l)}+b_m\ee^{-\ii\gamma_m(x-x_r)}\bigr)\ee^{\ii\beta_my}:
\end{equation}
$a$ enters through the left wall and $b$ through the right wall. $\mathcal
R[a,b]$ is a $\beta$-quasi-periodic solution of $(\Delta+\kappa_{\mathrm
e}^2)u=0$ in the whole plane; these are the plane-wave expansions at layer
interfaces of \cite[eq.~(3)]{BottenEtAl2004}, with $x$ and $y$ exchanged.

\emph{Cell field.} Let $G^{(\kappa)}(\mathbf z,\mathbf z')=\frac{\ii}{4}
H^{(1)}_0(\kappa|\mathbf z-\mathbf z'|)$, and let $\mathcal S^{(\kappa)}[\omega]$,
$\mathcal D^{(\kappa)}[\tau]$ be the single- and double-layer potentials on a
curve $\Gamma$. Consider a cell that contains one rod $\Omega$ with boundary
$\Gamma$, outward normal $\nu$, coefficient $\rho=\rho_{\mathrm c}$ in $\Omega$ and interior
wavenumber $\kappa_{\mathrm c}=\widehat k\rho_{\mathrm c}^{1/2}$. Given a density pair $\eta=(\tau;-\omega)$ on
$\Gamma$, proxy points $\mathbf p_i$ outside the closed cell with coefficients
$\alpha_i$, and amplitudes $a,b$, the \emph{cell field} is
\begin{equation}\label{eq:mfs-field}
  \widehat u=
  \begin{cases}
    \begin{aligned}[t]
      &\textstyle\sum_{n=-1}^{1}\ee^{\ii\beta nd}\bigl(\mathcal S^{(\kappa_{\mathrm e})}[\omega]
      +\mathcal D^{(\kappa_{\mathrm e})}[\tau]\bigr)(\cdot-nd\,e_y)\\
      &\quad\textstyle+\sum_i\alpha_iG^{(\kappa_{\mathrm e})}(\cdot,\mathbf p_i)+\mathcal R[a,b]
    \end{aligned}
    & \text{outside }\overline\Omega,\\[1ex]
    \mathcal S^{(\kappa_{\mathrm c})}[\omega]+\mathcal D^{(\kappa_{\mathrm c})}[\tau]
    & \text{in }\Omega .
  \end{cases}
\end{equation}
This is the quasi-periodic layer potential with its three nearest images kept
($e_y$ is the unit vector in $y$) and the remainder replaced by point sources: the quasi-periodic Green function
minus the three nearest images is a smooth Helmholtz solution near the closed
cell \cite{Linton1998}, the structure of
\cite{BarnettGreengard2010,BarnettGreengard2011,ChoBarnett2015}. The image
curves must have positive distance from the closed cell. If the cell has no rod,
the cell field is $\mathcal R[a,b]$ alone.

Three choices fix the field and are ours. (i) The density is the trigonometric
interpolant, in the parameter $t$ of the rod boundary $\boldsymbol r(t)$, of the
computed samples, and all layer potentials are exact integrals of
it; the Nystr\"om sums are not used as a field, since they are not $H^2$ up to
the interface. (ii) The proxy coefficients are fitted after the densities, by
least squares on collocation points (quasi-periodicity of value and $\partial_y$
on the top and bottom edges, outgoing Rayleigh waves on two auxiliary vertical
lines); unlike \cite{BarnettGreengard2010}, the densities are not re-solved
together with the proxies. (iii) In a lead cell the density and the proxy
coefficients are the cell solution for the incident field $\mathcal R[a,b]$,
hence linear in $(a;b)$, with one linear map per lead.

\subsection{Center and lead cells}\label{sec:center}
\emph{Center cell.} In $\mathcal C_0$, $\widehat u$ is the cell field
\eqref{eq:mfs-field} with rod $\Omega_0$ and Rayleigh term
\begin{equation}\label{eq:hom}
  \mathcal R[a_0,b_0]\quad\text{on }(x_-,x_+)\times I_d,
\end{equation}
the field entering through the two ports. It may be nonzero even if
$\Omega_0=\varnothing$, so quasi-periodic layer potentials alone do not
represent the center field. The ansatz follows the complementary-field
representation of \cite[App.~A]{BarnettGreengard2010} with the incident port
terms added; neither completeness nor uniqueness of this representation is used
below.

\emph{Lead cells.} In a half guide, the state of a wall is described by a vector
$v\in\CC^K$, and two matrices are given for each side. The matrix
$\mathsf P_\pm\in\CC^{K\times K}$ moves a state by one cell away from the center.
The matrix $\mathsf G_\pm\in\CC^{2K\times K}$ maps the state $v$ of the wall of
$\mathcal C_{0,\pm}$ nearer the center to the incident amplitudes $(a;b)$ of that
cell; $a$ is computed from $v$ and $b$ from $\mathsf P_\pm v$ on the right, and the
other way round on the left. Let $u_\pm[v]$ be the cell field of $\mathcal C_{0,\pm}$
with incident amplitudes $\mathsf G_\pm v$. With the lead states $c_\pm\in\CC^K$
at the ports, the field in the half guides is
\begin{equation}\label{eq:lead-field}
  \widehat u(\mathbf z\pm jL_\pm e_x)=u_\pm[\mathsf P_\pm^jc_\pm](\mathbf z),
  \qquad\mathbf z\in\mathcal C_{0,\pm},\ j\ge0 .
\end{equation}
The only property of $\mathsf P_\pm$ used below is $\|\mathsf P_\pm^p\|_2<1$ for some
$p\ge1$, which is checked (\cref{sec:computation}).

\emph{How the data are computed.} The densities come from a Nystr\"om
discretization of the transmission equations of each cell. For each lead, the
decaying states are those of the Bloch-mode eigenvalue problem for the cell
transfer matrix in the Rayleigh basis \cite[eq.~(7)]{BottenEtAl2004}, solved in
the equivalent pencil form built from the cell scattering matrix by an ordered
generalized Schur (QZ) decomposition. The columns of the resulting basis span
the decaying subspace, $\mathsf P_\pm$ is the one-cell translation in this basis
(nonnormal in general, not the diagonal of Bloch factors), and $\mathsf G_\pm$ is
read off from the basis. Finally $\widehat k$, $a_0$, $b_0$, $c_\pm$ solve the
matching of the center cell to the decaying subspaces of the two leads. For the
exact discrete lead problem neighboring cells agree on the modes $|m|\le M$; the
saved data do so only approximately, and the mismatch, the modes $|m|>M$, and
the discrete residuals appear as jumps on the shared walls.

In total, the candidate data are $\widehat k$, the center density, $a_0$, $b_0$
and proxy coefficients, the two linear maps of (iii), $\mathsf G_\pm$,
$\mathsf P_\pm$ and $c_\pm$. They are saved as binary64 numbers and treated as
exact; the field is homogeneous of degree one in the center data and $c_\pm$, so
the bound does not depend on their normalization.

\subsection{Properties of the trial field}\label{sec:field-properties}
The \emph{trial field} $\widehat u$ satisfies:
(P1) $(-\Delta-\mu\rho)\widehat u=0$ exactly in every homogeneous piece;
(P2) the seam mismatch between the top and bottom edges of each cell is an
interface jump;
(P3) for analytic rod boundaries, as in \cref{sec:results}, $\widehat u$ is $H^2$
up to the boundary of each piece and has ordinary one-sided traces on walls and
seams (\cref{lem:piecewise-h2});
(P4) $\widehat u\in\mathcal H$, which follows from the decay of $\mathsf P_\pm^j$
and is checked in \cref{sec:computation}.
It is not in $H^1_\beta(B)$.

\begin{lemma}[piecewise regularity of the trial field]\label{lem:piecewise-h2}
Assume that every material boundary has a regular embedded analytic periodic
parametrization, that the image curves and proxies are at positive distance from
the closed cells, that the densities are trigonometric polynomials in that
parameter, and that \eqref{eq:mfs-field} uses one $\kappa_{\mathrm e}$ in the
background. Then $\widehat u$ is $H^2$ up to the boundary in each homogeneous
piece and has ordinary value and normal traces on both sides of every wall and
seam.
\end{lemma}

\begin{proof}
Far images, proxies and plane waves are smooth near the closed pieces, so it
suffices to consider $v=\mathcal S^{(\kappa)}[\omega]+\mathcal D^{(\kappa)}[\tau]$ on
one curve $\Gamma$. In the sense of distributions $(\Delta+\kappa^2)v=
-\omega\delta_\Gamma+\operatorname{div}(\tau\nu\delta_\Gamma)$. In normal coordinates
let $w=0$ inside and $w=(\tau(s)-t\omega(s))\chi(t)$ outside, with a smooth cutoff
$\chi$. Then $w$ has value jump $\tau$ and normal jump $-\omega$, so
$(\Delta+\kappa^2)w=F-\omega\delta_\Gamma+\operatorname{div}(\tau\nu\delta_\Gamma)$ with
$F\in L^2_{\mathrm{loc}}$. Hence $(\Delta+\kappa^2)(v-w)=-F$ and $v-w\in
H^2_{\mathrm{loc}}$, while $w$ is smooth up to $\Gamma$ from either side. Near walls
and seams the one-sided cell formulas extend smoothly across, so there are no
corner singularities.
\end{proof}
